\section{Related Works}

\subsection{State Space Models}

Recently, the State Space Models (SSMs) have shown significant effectiveness of state space transformation in capturing the dynamics and dependencies of language sequences.
~\cite{ssms4} introduces a structured state-space sequence model (S4), specifically designed to model long-range dependencies, boasting the advantage of linear complexity. 
Based on it,
various models have been developed
(\textit{e.g.}, S5~\cite{ssms5}, H3~\cite{h3layer} and GSS~\cite{mehta2022long}),
% S6, mamba
and Mamba~\cite{gu2023mamba} distinguishes itself 
by introducing a data-dependent SSM layer and a selection mechanism using parallel scan~(S6).
Compared to transformers~\cite{lu2019vilbert,gpt3} based on quadratic-complexity attention, 
Mamba excels at processing long sequences with linear complexity.

In the vision domain,
\cite{ssms4} first applies SSM in pixel-level image classification,
and \cite{vis4mer} uses S4 to handle the long-range temporal dependencies for movie clip classification.
Besides,
the great potential of Mamba motivates a series of works~\cite{vim,vmamba,liang2024pointmamba,vivim,guo2024mambair,he2024pan,wang2024graph},
which demonstrates Mamba's better performances and higher GPU efficiency than Transformer on visual downstream tasks like object detection and semantic segmentation. 
Different from the previous works,
our \ModelName\ is a purely SSM-based video model,
showcasing great efficiency and effectiveness for both short-term and long-term video understanding.


\subsection{Video Understanding}

Video understanding stands as a cornerstone in the domain of computer vision, 
whose significance is further amplified by the burgeoning growth of short video platforms. 
To bolster this field, 
numerous datasets equipped with extensive data and meticulous human annotations have been developed, aiming to enhance human action recognition capabilities. 
Notable examples include UCF101~\cite{ucf101} and Kinetics dataset~\cite{k400,k600,k700}, 
which have played pivotal roles in benchmarking progress. 
Furthermore, other datasets~\cite{activitynet,thumos,charades_sta,ava,finegym,liu2022fineaction} provide annotated activity videos tailored for action localization,
fostering deeper research into human activities.
Beyond action recognition, 
the advent of large-scale video-text datasets~\cite{msrvtt,msvd,activitynet_qa,youcook2, miech2019howto100m,internvid} extends the utility of video understanding into the realm of multi-modality tasks,
such as video captioning, retrieval and question answering,
thereby broadening the application spectrum. 

As for the architecture,
it has evolved from using CNN which extracts features from video frames, to more advanced techniques.
Initially, 3D CNNs~\cite{c3d,i3d,x3d,r(2+1)d} expanded the traditional 2D CNN architecture to capture videos' spatio-temporal information. 
Two-Stream~\cite{two_stream}, which combines spatial and temporal streams, 
TSN~\cite{tsn}, which proposes sparse sampling,
and SlowFast~\cite{slowfast}, which uses parallel networks to capture semantics and rapid movements, further enhance action recognition capacity. 
The introduction of attention-based models~\cite{timesformer,vivit,motionformer,stam,morphmlp}, like TimeSformer~\cite{timesformer} and ViViT~\cite{vivit}, 
marked a significant advancement by effectively capturing long-range dependencies within video sequences, 
enhancing temporal relationship understanding. Recent developments~\cite{uniformer,uniformerv2,Wang2022InternVideoGV,video_swin} have focused on accurate video transformer, with innovations like the VideoSwin's window attention~\cite{video_swin} and the UniFormer's integration of convolution and self-attention mechanisms~\cite{uniformer}, 
aiming to balance computational efficiency with performance. 
Despite these models' achievements in various tasks, 
they often come with high computational costs for long sequences. 
In contrast,
our \ModelName\  introduces a linear-complexity operator for efficient long-term modeling, outperforming existing methods with faster speed and lower GPU consumption.
